% !PRINTED 11
\weberpage{11}

The natural numbers themselves are inserted into the manifold of
rational numbers if by the natural number $m$ one understands the
ratio $m:1$.

We now return to any measurable manifold $\mathfrak M$ and take from
it any two elements $a$ and $b$. If one chooses, which is always
possible, two natural numbers $m,n$ so that $na>mb$, then the ratio
$a:b$ is called greater than the rational ratio $m:n$, or
\[
 {a\over b}>{m\over n},
\]
and if $m:n>p:q$, then also $a:b>p:q$. Similarly, if $n'a<m'b$, then
\[
 {a\over b}<{m'\over n'}.
\]

If $a:b>m:n$, then one can insert one, and consequently also
arbitrarily many, rational numbers $m_1:n_1$ such that
\[
 {a\over b}>{m_1\over n_1}>{m\over n}
\]
hold. That is, one can insert arbitrarily many rational ratios
\gesperrt{between} $a:b$ and $m:n$. To show this, choose an arbitrary
whole number $k$ and determine the whole number $h$ so that
$h(na-mb)>kb$, which is always possible. Then
\[
 {a\over b}>{hm+k\over hn}>{m\over n}.
\]
And in exactly the same way, if $n'a<m'b$ and
$h'(m'b-n'a)>k'a$, one obtains
\[
 {a\over b}<{h'm'\over h'n'+k'}<{m'\over n'}.
\]

If now $a:b$ and $\alpha:\beta$ are any two ratios, which for brevity
we shall also denote by $e,\epsilon$, and whose elements belong to the
same or to different manifolds, then two cases are possible: 1) no
rational ratio $\mu$ lies between $e$ and $\epsilon$, or 2) a rational
ratio lies between $e$ and $\epsilon$.

In case 1) the two ratios $e$ and $\epsilon$ are called
\gesperrt{equal to one another}. One sees that two ratios which are
equal to a third are also equal to one another; for if
$e<\mu<\epsilon$, then every other ratio $e'$ is either smaller
%CONT>>
